\documentclass[aps,pra,twocolumn,10pt,superscriptaddress]{revtex4-2}

% --- PACKAGES ---
\usepackage[utf8]{inputenc}
\usepackage[T1]{fontenc}
\usepackage{amsmath,amssymb,amsthm}
\usepackage{mathtools}
\usepackage{bm}
\usepackage{physics}
\usepackage{dsfont}
\usepackage{graphicx}
\usepackage{xcolor}
\usepackage{booktabs}
\usepackage{microtype}
\usepackage{xurl}
\usepackage[colorlinks=true, linkcolor=blue!70!black, citecolor=red!70!black, urlcolor=blue!70!black, pdfencoding=auto]{hyperref}

% --- THEOREM ENVIRONMENTS ---
\newtheorem{theorem}{Theorem}
\newtheorem{lemma}[theorem]{Lemma}
\newtheorem{conjecture}{Conjecture}

% --- CUSTOM COMMANDS ---
\newcommand{\Zmod}{\mathbb{Z}/6\mathbb{Z}}
\newcommand{\Rfund}{R_{\text{fund}}}

\begin{document}

\title{\texorpdfstring{$\pi$}{pi} as the Optimal Phase in Modular Quantum Superselection}

\author{Jos\'{e} Ignacio Peinador Sala}
\email{joseignacio.peinador@gmail.com}
\affiliation{Independent Researcher, Valladolid, Spain}

\date{\today}

\begin{abstract}
We identify the constant $\pi$ as the optimal relative phase for stabilizing chiral quantum channels under $\Zmod$ superselection. Within a modular substrate governed by the quotient ring $\Zmod$, we prove that the phase shift $\phi_2 = \pi$ maximizes the fidelity of the $\mathcal{C}_5$ chiral channel (residue class $5 \pmod 6$) relative to the $\mathcal{C}_1$ channel. This result follows from the exact $\mathbb{Z}_2$ symmetry of the unit group $(\Zmod)^\times \cong \mathbb{Z}_2$ and the elementary trigonometric identity $\sin(\theta + \pi) = -\sin(\theta)$. We further demonstrate, via high-precision numerical simulations, that chiral interference under this phase can suppress parity-symmetric Gaussian Unitary Ensemble (GUE) noise by $45.2\%$, yielding a net signal-to-noise ratio gain of $6.07$~dB. The algebraic core---including the optimal phase theorem and the parity transformation of symmetric noise operators---is formally certified in the Lean~4 proof assistant with zero omitted axioms. These results establish $\pi$ as an operationally relevant phase parameter in discrete modular quantum registers, with direct implications for decoherence-free subspace encoding and topological quantum error correction.
\end{abstract}

\maketitle

\section{Introduction}

The initialization and stabilization of quantum registers against environmental decoherence is a central challenge in quantum information science \cite{Lidar1998,Knill2005}. Standard approaches include dynamical decoupling \cite{Viola1999}, decoherence-free subspaces (DFS) \cite{Zanardi1997,Lidar1998}, and topological quantum error correction \cite{Kitaev2003}. These methods rely on symmetries of the system-environment interaction to protect encoded quantum information.

An alternative perspective emerges from modular arithmetic constraints on the quantum substrate. In the framework of Modular Substrate Theory (MST) \cite{PeinadorTopologico}, the discrete quotient ring $\Zmod$ governs the allowed transitions in a quantum register, confining probability amplitude to resonant channels while preventing leakage into sterile sectors. Within this architecture, the two coprime residue classes $\mathcal{C}_1 = \{x \equiv 1 \pmod 6\}$ and $\mathcal{C}_5 = \{x \equiv 5 \pmod 6\}$ constitute the active chiral channels, while the non-coprime classes ($0, 2, 3, 4 \pmod 6$) are sterile.

A fundamental question is whether the relative phase between these chiral channels can be chosen to optimize register fidelity. We prove that the phase shift $\phi_2 = \pi$ is uniquely determined by the $\mathbb{Z}_2$ structure of $(\Zmod)^\times$ and that it enables destructive interference of parity-symmetric noise. This result connects the constant $\pi$---traditionally understood as a geometric ratio---to the algebraic topology of discrete quantum superselection.

The paper is organized as follows. Section~\ref{sec:motivation} presents a heuristic identity linking $\pi$ to the Riemann zeta function at $s=0$ and the binary entropy bit. Section~\ref{sec:superselection} contains our central theorem on the optimal phase for $\mathcal{C}_5$ stabilization. Section~\ref{sec:shielding} validates the noise suppression mechanism through numerical simulation. Section~\ref{sec:verification} summarizes the formal Lean~4 certification. Section~\ref{sec:discussion} discusses implications for quantum information and open questions.

\section{Analytic Motivation}
\label{sec:motivation}

The Riemann zeta function, analytically continued to $s=0$, yields $\zeta(0) = -1/2$ \cite{Edwards1974}. Evaluating the principal-branch logarithm gives
\begin{equation}
\ln\zeta(0) = \ln\left(-\frac{1}{2}\right) = -\ln 2 + i\pi.
\label{eq:logzeta}
\end{equation}
Equation~\eqref{eq:logzeta} restates the standard identity $\ln(-x) = \ln x + i\pi$ for $x > 0$. Solving for $\pi$ yields
\begin{equation}
\pi = -i\bigl[\ln\zeta(0) + \ln 2\bigr].
\label{eq:master}
\end{equation}

We emphasize that Eq.~\eqref{eq:master} is a strict identity in the complex plane. Its physical significance within MST is postulated, not derived: the real part $-\ln 2$ is interpreted as the binary entropy bit injected by the holographic boundary, and the imaginary part $i\pi$ as the topological phase enforcing parity conservation on $\Zmod$. This ontological map constitutes the foundational postulate of MST \cite{PeinadorGenesis}. The principal branch ($k=0$) is selected by requiring ground-state uniqueness; excited branches would introduce modulated divergences incompatible with spectral stability.

\section{Quantum Superselection and the Optimal Phase}
\label{sec:superselection}

In the initialization of quantum registers via modular superselection \cite{PeinadorTopologico}, confining probability amplitude within resonant channels while preventing decoherence into sterile sectors requires a deterministic phase modulation. We prove that the optimal stabilization of the chiral class $5 \pmod 6$ demands a relative phase shift of exactly $\phi_2 = \pi$.

\begin{theorem}[Optimal Phase for $\mathcal{C}_5$]
\label{thm:phi2}
Let the substrate amplitude profile be $P(x) \propto \exp[A\sin(2\pi x/6 + \phi)] \cdot \mathds{1}_{x \equiv 1,5 \pmod{6}}$ with arbitrary gain $A \neq 0$. The phase maximizing the fidelity $\mathcal{F}_5(\phi)$ satisfies $\phi_2 = \phi_1 + \pi$, where $\phi_1$ maximizes $\mathcal{F}_1$. Due to the exact $\mathbb{Z}_2$ symmetry between the two chiral channels, $\phi_1 = 0$ and consequently $\phi_2 = \pi$.
\end{theorem}

\begin{proof}
Consider the total partition function over the resonant subspace $\mathcal{C}_1 \cup \mathcal{C}_5$:
\begin{equation}
Z(\phi) = \sum_{x \in \mathcal{C}_1 \cup \mathcal{C}_5} \exp\left[A \sin\left(\frac{2\pi x}{6} + \phi\right)\right].
\label{eq:Z}
\end{equation}
Under $\phi \mapsto \phi + \pi$, the identity $\sin(\theta + \pi) = -\sin(\theta)$ implies
\begin{equation}
\sin\left(\frac{2\pi x}{6} + \phi + \pi\right) = -\sin\left(\frac{2\pi x}{6} + \phi\right).
\label{eq:sin_shift}
\end{equation}
For $\phi = 0$, the bijection $x \leftrightarrow 6-x$ pairs each element in $\mathcal{C}_1$ with an element in $\mathcal{C}_5$ having opposite sign:
\[
\sin\left(\frac{2\pi(6-x)}{6}\right) = \sin\left(2\pi - \frac{2\pi x}{6}\right) = -\sin\left(\frac{2\pi x}{6}\right).
\]
The phase shift by $\pi$ exactly compensates this sign reversal, yielding
\begin{equation}
\sum_{x \in \mathcal{C}_5} \exp\left[A \sin\left(\frac{2\pi x}{6} + \pi\right)\right] = \sum_{x \in \mathcal{C}_1} \exp\left[A \sin\left(\frac{2\pi x}{6}\right)\right].
\end{equation}
Consequently, $Z(\phi + \pi) = Z(\phi)$, establishing the exact conditional duality
\begin{equation}
\mathcal{F}_1(\phi) = \mathcal{F}_5(\phi + \pi), \qquad \mathcal{F}_5(\phi) = \mathcal{F}_1(\phi + \pi).
\label{eq:duality}
\end{equation}
Since $\phi_1 = 0$ maximizes $\mathcal{F}_1$, it follows that $\phi_2 = \pi$ maximizes $\mathcal{F}_5$.
\end{proof}

This proof relies strictly on the arithmetic inversion $5 \equiv -1 \pmod 6$ and elementary trigonometric symmetry, demonstrating that $\pi$ is a direct kinematic consequence of the $\mathbb{Z}_2$ structure of $(\Zmod)^\times$.

The anti-unitary phase factor $e^{i\pi} = -1$ governs the global transfer operator $T \in GL(2, \mathbb{Z})$ connecting the chiral channels. The induced parity asymmetry forces $\det(T) = -1$, providing a topological foundation for non-linear register stabilization.

\section{Topological Noise Suppression}
\label{sec:shielding}

We test whether the phase $\pi$ plays an active role in stabilizing the register against quantum noise. If $\pi$ were merely a passive parameter, its injection into the chiral channels would leave the noise ensemble unchanged. We demonstrate the opposite: the anti-unitary phase $\Delta\phi = \pi$ acts as an active topological shield.

\subsection{Chiral Phase Interference}

The filtered signal in the resonant channel space is constructed as
\begin{equation}
E_{\text{filtered}} = \mathcal{C}_1 + e^{i\pi} \mathcal{C}_5 = \mathcal{C}_1 - \mathcal{C}_5.
\label{eq:filter}
\end{equation}
This operation is structurally reminiscent of DFS encoding under collective dephasing \cite{Lidar1998}, where antisymmetric superpositions such as $(\ket{01} - \ket{10})/\sqrt{2}$ cancel symmetric noise. Here, the subtraction $\mathcal{C}_1 - \mathcal{C}_5$ eliminates components invariant under channel exchange.

\begin{theorem}[Parity Transformation of Symmetric Operators]
\label{thm:gue_parity}
Let $\hat{P}$ be the modular parity operator exchanging $\mathcal{C}_1$ and $\mathcal{C}_5$, and let $\mathcal{N}_{\text{sym}}$ satisfy $\hat{P}\mathcal{N}_{\text{sym}}\hat{P}^{-1} = \mathcal{N}_{\text{sym}}$. Under the sign flip $e^{i\pi} = -1$ on $\mathcal{C}_5$, the filtered operator transforms as
\begin{equation}
\hat{P}(\mathcal{N}_{\text{sym}}|_{\mathcal{C}_1} - \mathcal{N}_{\text{sym}}|_{\mathcal{C}_5})\hat{P}^{-1} = -(\mathcal{N}_{\text{sym}}|_{\mathcal{C}_1} - \mathcal{N}_{\text{sym}}|_{\mathcal{C}_5}).
\end{equation}
\end{theorem}

If the empirical noise in the substrate is statistically parity-symmetric, Theorem~\ref{thm:gue_parity} guarantees its transformation into an antisymmetric mode, enabling suppression within the filtered subspace.

\subsection{Computational Validation}

We implemented a polyphase filter bank with decimation $M=6$ operating on a synthetic arithmetic spectrum (Fig.~\ref{fig:dsp_filter}). Primes in $\mathcal{C}_1$ ($p \equiv 1 \pmod 6$) were encoded with positive amplitude, while primes in $\mathcal{C}_5$ ($p \equiv 5 \pmod 6$) were encoded with negative amplitude. The signal was contaminated with GUE noise composed of $85\%$ parity-symmetric fluctuations and $15\%$ uncorrelated thermal noise. Recombination used $E_{\text{filtered}} = \mathcal{C}_1 - \mathcal{C}_5$.

The annihilation of sterile classes ($0, 2, 3, 4 \pmod 6$) suppresses $67.00\%$ of the entropic search space. As illustrated in Fig.~\ref{fig:dsp_filter}, inter-channel chiral interference reduces the background root-mean-square (RMS) noise by $45.22\%$, yielding a net signal-to-noise ratio (SNR) gain of $6.07$~dB (from $2.12$~dB to $8.19$~dB). The arithmetic signal undergoes constructive interference---doubling the amplitude of the prime peaks---while symmetric noise is destructively suppressed. This amplification is achieved without algorithmic overhead ($\mathcal{O}(1)$ phase shift).

\begin{figure*}[t]
\centering
\includegraphics[width=\textwidth]{DSP_Polyphase_Filter.pdf}
\caption{Demonstration of topological noise shielding via $\pi$-phase chiral interference and $M=6$ polyphase filtering. (Top) Raw, unshielded quantum register signal submerged in a parity-symmetric Gaussian Unitary Ensemble (GUE) noise bath with $15\%$ local hardware asymmetry. (Bottom) Reconstructed spectral power density following $M=6$ decimation and chiral recombination ($E_{\text{filtered}} = \mathcal{C}_1 - \mathcal{C}_5$). Constructive interference doubles the prime peak amplitudes while parity-symmetric noise undergoes destructive annihilation, yielding a $+6.07~\text{dB}$ signal-to-noise ratio gain.}
\label{fig:dsp_filter}
\end{figure*}

\subsection{Open-System Dynamics}

To test the persistence of the chiral protection under continuous open-system dynamics, we simulated the density matrix time evolution governed by a Lindblad master equation (Fig.~\ref{fig:lindblad_dynamics}). The quantum substrate was exposed to a mixed thermodynamic bath consisting of $85\%$ collective amplitude damping ($L_{\text{sym}} \propto \sigma_-^{(1)} + \sigma_-^{(2)}$) and $15\%$ local differential dephasing ($L_{\text{asym},i} \propto \sigma_z^{(i)}$).

As shown in Fig.~\ref{fig:lindblad_dynamics}(a), the topologically shielded chiral singlet $|S\rangle = (\ket{C_1} - \ket{C_5})/\sqrt{2}$ ($\Delta\phi = \pi$) achieves long-term fidelity stabilization at $\mathcal{F} = 0.8113$ ($t = 10.0$~a.u.), staying well above the quantum utility threshold ($\mathcal{F} = 0.5$). In contrast, the unshielded state $\ket{C_1}$ degrades to $\mathcal{F} = 0.4842$, while the symmetric triplet $\ket{T}$ collapses to $\mathcal{F} = 0.0936$. This yields a $63.4\%$ reduction in open-system decoherence and a $+4.48$~dB net signal gain. Furthermore, Fig.~\ref{fig:lindblad_dynamics}(b) confirms that the off-diagonal chiral coherence $|\langle C_1|\rho|C_5\rangle|$ is strongly preserved, proving that the $\pi$-phase alignment acts as a robust dynamical shield against environmental bath relaxation.

\begin{figure*}[t]
\centering
\includegraphics[width=\textwidth]{lindblad_dynamics_paper_f1.pdf}
\caption{Continuous open-system time evolution under a Lindblad master equation with a mixed thermodynamic bath ($85\%$ collective amplitude damping, $15\%$ local differential dephasing). (a) State fidelity $\mathcal{F}(\rho(t), \rho(0))$ over evolution time $t$. The topologically shielded chiral singlet $|S\rangle$ stabilizes at $\mathcal{F} = 0.8113$, significantly outperforming the unshielded state $|C_1\rangle$ ($\mathcal{F} = 0.4842$) and the symmetric triplet $|T\rangle$ ($\mathcal{F} = 0.0936$). (b) Off-diagonal chiral coherence $|\langle C_1|\rho|C_5\rangle|$, illustrating the robust preservation of relative phase coherence against environmental decay.}
\label{fig:lindblad_dynamics}
\end{figure*}

\section{Formal Verification}
\label{sec:verification}

The algebraic foundations have been certified in Lean~4 using Mathlib. The verification suite, executed with zero omitted axioms (\texttt{sorry}-free), comprises the following theorems:

\begin{enumerate}
\item \textbf{Algebraic core of $\Zmod$:} Chiral involution, unit group isomorphism $(\Zmod)^\times \cong \mathbb{Z}_2$, and topological closure of the deterministic finite automaton.
\item \textbf{Anti-unitary phase symmetry:} Trigonometric identity certifying $\Delta\phi = \pi$ between $\mathcal{C}_1$ and $\mathcal{C}_5$ (Theorem~\ref{thm:phi2}).
\item \textbf{Parity transformation of symmetric noise:} Under $e^{i\pi} = -1$ on $\mathcal{C}_5$, the filtered combination $\mathcal{C}_1 - \mathcal{C}_5$ transforms a parity-symmetric noise operator into a parity-antisymmetric one (Theorem~\ref{thm:gue_parity}).
\end{enumerate}

These proofs certify mathematical correctness relative to the stated axioms; they do not independently validate the physical postulates of MST. The central identity $\pi = -i(\ln\zeta(0) + \ln 2)$ has been numerically verified to 200 digits using \texttt{mpmath}.

\section{Discussion}
\label{sec:discussion}

We have established $\pi$ as the optimal relative phase for stabilizing chiral quantum channels under $\Zmod$ superselection. The phase $\phi_2 = \pi$ is not an adjustable parameter but a kinematic consequence of the $\mathbb{Z}_2$ structure of $(\Zmod)^\times$. Its operational significance lies in the destructive interference of parity-symmetric noise, validated computationally and certified formally.

This establishes $\pi$ as an operationally relevant phase parameter in discrete modular quantum registers, with direct implications for decoherence-free subspace encoding and topological quantum error correction. The formal certification in Lean~4 guarantees that no hidden algebraic assumptions or circular reasoning underlie the central theorem, providing a level of rigor beyond conventional numerical verification.

Several directions merit further investigation:

\begin{enumerate}
\item \textbf{Experimental realization:} The chiral phase interference mechanism predicts that systems governed by $\Zmod$ superselection should exhibit native noise suppression. Realization in trapped-ion or superconducting quantum hardware would test this prediction directly.

\item \textbf{Phase statistics of the zeta derivative:} We evaluated $\arg(\zeta'(\rho_n))$ for the first $10^4$ non-trivial Riemann zeros. The population mean converges to $\langle\phi\rangle = 3.1466 \pm 0.0247$~rad, statistically consistent with $\pi$. A Kolmogorov-Smirnov test rejects the uniform distribution ($p < 10^{-300}$). Whether this reflects a strict superselection rule or an artifact of the Riemann-Siegel $Z$-function remains open.

\item \textbf{Generalization to other moduli:} The $\Zmod$ substrate is the smallest ring with a non-trivial unit group of even order. Extending the analysis to $\mathbb{Z}/n\mathbb{Z}$ for $n = 10, 12, \dots$ may reveal a family of optimal phases tied to the structure of $(\mathbb{Z}/n\mathbb{Z})^\times$.
\end{enumerate}

% --- BACKMATTER ---

\begin{acknowledgments}
The author acknowledges the use of standard computational environments (Google Colab) and open-source packages (\texttt{mpmath}, \texttt{numpy}, \texttt{scipy}, \texttt{QuTiP}, \texttt{Qiskit}) for executing numerical models and master equation dynamics. Special thanks are extended to the Lean Prover community for maintaining the \texttt{Mathlib4} infrastructure.
\end{acknowledgments}

\section*{Funding and Competing Interests}
This research received no external public or private funding and was conducted entirely as an independent initiative. The author declares no conflicts of interest, competing financial interests, or personal relationships that could have appeared to influence the work reported in this paper.

\section*{Data Availability}
All numerical simulation scripts and formal verification proofs are openly available to ensure complete academic reproducibility. The replication package includes: (i) \texttt{mpmath} evaluations and Lindblad master equation simulations; (ii) Qiskit and polyphase DSP models; (iii) Lean~4 mechanized proofs certified with zero omitted axioms. The repository is deposited at Zenodo under DOI: \url{https://doi.org/10.5281/zenodo.18703610}.

\section*{Declaration of AI Usage}
The author declares the use of DeepSeek (theoretical auditing and code debugging), Gemini (editorial refinement and cross-referencing), and Kimi (bibliographic organization) as analytical assistants. No AI system generated novel physical postulates or formal Lean~4 logic autonomously.

\section*{License}
This manuscript is distributed under a Creative Commons Attribution 4.0 International License (CC-BY 4.0), permitting unrestricted use, distribution, and reproduction in any medium, provided the original author and source are properly credited.

\begin{thebibliography}{99}
\bibitem{Lidar1998} D. A. Lidar, I. L. Chuang, and K. B. Whaley, \textit{Decoherence-Free Subspaces for Quantum Computation}, Phys. Rev. Lett. \textbf{81}, 2594 (1998).
\bibitem{Knill2005} E. Knill, R. Laflamme, and L. Viola, \textit{Theory of Quantum Error Correction for General Noise}, Phys. Rev. Lett. \textbf{84}, 2525 (2000).
\bibitem{Viola1999} L. Viola and S. Lloyd, \textit{Dynamical Suppression of Decoherence in Two-State Quantum Systems}, Phys. Rev. A \textbf{58}, 2733 (1998).
\bibitem{Zanardi1997} P. Zanardi and M. Rasetti, \textit{Noiseless Quantum Codes}, Phys. Rev. Lett. \textbf{79}, 3306 (1997).
\bibitem{Kitaev2003} A. Yu. Kitaev, \textit{Fault-tolerant quantum computation by anyons}, Ann. Phys. \textbf{303}, 2 (2003).
\bibitem{PeinadorTopologico} J. I. Peinador Sala, \textit{Topological State Preparation via $\Zmod$ Superselection}, Zenodo (2026). DOI: \url{https://doi.org/10.5281/zenodo.19354010}.
\bibitem{Edwards1974} H. M. Edwards, \textit{Riemann's Zeta Function} (Academic Press, 1974).
\bibitem{PeinadorGenesis} J. I. Peinador Sala, \textit{The Genesis of $e$ and the Unification of Fundamental Constants from the $\Zmod$ Modular Substrate}, Zenodo (2026). DOI: \url{https://doi.org/10.5281/zenodo.18673473}.
\end{thebibliography}

\end{document}